\documentclass[11pt]{amsart}

% Encoding and fonts
\usepackage[T1]{fontenc}
\usepackage[utf8]{inputenc}
\usepackage{lmodern}

% Math packages
\usepackage{amsmath}
\usepackage{amssymb}
\usepackage{amsfonts}
\usepackage{amsthm}
\usepackage{mathtools}
\usepackage{mathrsfs}

% Layout and utilities
\usepackage{geometry}
\usepackage{enumitem}

\geometry{margin=1in}

\numberwithin{equation}{section}
\setcounter{tocdepth}{3}
\setcounter{secnumdepth}{3}

% Theorem-like environments
\theoremstyle{plain}
\newtheorem{theorem}{Theorem}[section]
\newtheorem{proposition}[theorem]{Proposition}
\newtheorem{lemma}[theorem]{Lemma}
\newtheorem{corollary}[theorem]{Corollary}
\newtheorem{conjecture}[theorem]{Conjecture}

\theoremstyle{definition}
\newtheorem{definition}[theorem]{Definition}
\newtheorem{experiment}[theorem]{Experiment}
\newtheorem{heuristic}[theorem]{Heuristic}
\newtheorem{question}[theorem]{Question}

\theoremstyle{remark}
\newtheorem{remark}[theorem]{Remark}
\newtheorem{warning}[theorem]{Warning}

\newcommand{\N}{\mathbb N}
\newcommand{\C}{\mathbb C}
\newcommand{\R}{\mathbb R}
\newcommand{\calP}{\mathcal P}
\newcommand{\CMZ}{\mathrm{CMZ}}
\newcommand{\RH}{\mathrm{RH}}
\newcommand{\ord}{\operatorname{ord}}
\newcommand{\Res}{\operatorname*{Res}}
\newcommand{\sgn}{\operatorname{sgn}}
\newcommand{\eps}{\varepsilon}

% Hyperref is loaded last.
\usepackage[hidelinks]{hyperref}

\title{An Ordinal-Prime CMZ Deformation Criterion\\
for the Riemann Hypothesis (v3)}
\author{@kasnyscience}
\date{}

\begin{document}

\begin{abstract}
We give a self-contained arithmetic-to-dynamical bridge between the Riemann
hypothesis and a two-channel logarithmic stretching observable motivated by
the multiscale mechanisms of C\'ordoba, Mart\'inez-Zoroa, and Zheng.  The
arithmetic input is the sign-specific ordinal-prime transform $S$, its
summatory function $SS$, and the Dirichlet series $\alpha$.  An exact finite
Abel identity, including its unfinished-pair boundary term, identifies the
residual Dirichlet polynomial $\delta_Q(\sigma)$ with the logarithm of the
relative deformation of two positive CMZ-type channels.  We prove that the
Riemann hypothesis is equivalent to tail stability of this relative
deformation for every real $\sigma>3/2$.  We also record the domains of
absolute convergence, all justified interchanges, the meromorphic blocks of
$\alpha$, and the singularities that cannot be transferred to the full
series without additional continuation.  The equivalence is rigorous; a
realisation by the published Euler gluing scheme is stated separately as an
open PDE problem and is not asserted here.  Three concrete attempts to feed
the ordinal-prime data into a genuine forced Euler or IPM tower are examined:
each is shown either to require an extension of the cited theorems that is
not on record, or to encode the arithmetic input without constraining it.  A
bounded-ratio generalisation of the unforced layered self-interaction bound
of C\'ordoba, Mart\'inez-Zoroa, and Zheng is proved along the way.  The
precise sense in which completing any such realisation would still not
decide the Riemann hypothesis is stated as
Warning~\ref{warn:no-progress-without-independent-bound}.  Two further
candidates are then excluded outright rather than left unproved: reading
the data off a single tower's amplitude ratio conflicts with an
already-proved sign-definite lemma, and reading it off a two-tower
comparison of activated packets conflicts, for large arithmetic index $q$,
with an unconditional quadratic-summability bound proved here for the
target sequence.  The one candidate excluded by neither obstruction ---
reading the data off a residual term the construction already produces,
rather than off an amplified response --- is posed as an explicit open
question.
\end{abstract}

\maketitle

\tableofcontents

\bigskip

\begin{remark}[Mathematical status and provenance]
\label{rem:status}
Statements marked \textup{[LIT]} are imported from the cited literature.
Statements marked \textup{[DED]} are proved in this article from the stated
definitions.  Statements marked \textup{[ADAPT]} adapt a published mechanism
to a different object.  A statement marked \textup{[OPEN]} is not proved.
The exact ordinal identities and the CMZ--RH equivalence are \textup{[DED]}.
The scalar stretching law is an \textup{[ADAPT]} of the CMZ amplitude
mechanism.  No theorem in this article claims that the existing Euler,
Boussinesq, or IPM constructions supply the missing arithmetic tail estimate.
\end{remark}

\section{Introduction}
\label{sec:introduction}

For $s=\sigma+it$ with $\sigma>1$, the Riemann zeta function is defined by
the absolutely convergent Dirichlet series and Euler product
\begin{equation}
\zeta(s)=\sum_{n=1}^{\infty}\frac1{n^s}
=\prod_p\left(1-p^{-s}\right)^{-1}.
\label{eq:zeta-series-product}
\end{equation}
It extends meromorphically to $\C$, has a simple pole at $s=1$, and satisfies
a functional equation relating $s$ and $1-s$; see
\cite{Titchmarsh1986,Apostol1976}.  Its nontrivial zeros lie in
$0<\Re s<1$.  The Riemann hypothesis is the assertion
\begin{equation}
\RH:\qquad \zeta(\rho)=0,\quad 0<\Re\rho<1
\quad\Longrightarrow\quad \Re\rho=\frac12.
\label{eq:RH-statement}
\end{equation}

The bridge developed here begins with four classical arithmetic functions.
The M\"obius function $\mu$, the Liouville function $\lambda$, and the von
Mangoldt function $\Lambda$ are
\begin{align}
\mu(n)&=
\begin{cases}
1,&n=1,\\
(-1)^r,&n\text{ is a product of }r\text{ distinct primes},\\
0,&p^2\mid n\text{ for some prime }p,
\end{cases}
\label{eq:mobius-definition}\\
\lambda(n)&=(-1)^{\Omega(n)},
\label{eq:liouville-definition}\\
\Lambda(n)&=
\begin{cases}
\log p,&n=p^k\text{ with }p\text{ prime and }k\ge1,\\
0,&\text{otherwise}.
\end{cases}
\label{eq:mangoldt-definition}
\end{align}
Here $\Omega(n)$ counts prime factors with multiplicity.  Absolute convergence
and multiplication of Euler factors give, for $\sigma>1$,
\begin{align}
\frac1{\zeta(s)}&=\sum_{n=1}^{\infty}\frac{\mu(n)}{n^s},
\label{eq:mobius-zeta}\\
\frac{\zeta(2s)}{\zeta(s)}&=
\sum_{n=1}^{\infty}\frac{\lambda(n)}{n^s},
\label{eq:liouville-zeta}\\
-\frac{\zeta'(s)}{\zeta(s)}&=
\sum_{n=1}^{\infty}\frac{\Lambda(n)}{n^s}.
\label{eq:mangoldt-zeta}
\end{align}
The first identity makes cancellation in $\mu$ directly sensitive to zeros of
$\zeta$.  The second and third identities explain why Liouville signs and
prime-power weights are natural comparison objects.  The construction below
uses $\mu$ itself; $\lambda$ and $\Lambda$ are introductory reference points,
not substituted coefficients.

Write
\begin{equation}
M(x)=\sum_{n\le x}\mu(n).
\label{eq:mertens-definition}
\end{equation}
The classical Littlewood criterion states
\begin{equation}
\RH\quad\Longleftrightarrow\quad
M(x)=O_\eta\!\left(x^{1/2+\eta}\right)
\quad\text{for every }\eta>0.
\label{eq:littlewood-criterion}
\end{equation}
We use this imported theorem in its stated pointwise form
\cite{Littlewood1912,Titchmarsh1986}.  The objective of this article is not to
replace \eqref{eq:littlewood-criterion} by an unproved estimate under a new
name.  It is to identify exactly how an ordinal-prime observable enters a CMZ
logarithmic stretching law and to state the resulting equivalence without
discarding cut-off or analytic-domain information.

\section{The ordinal-prime transform}
\label{sec:ordinal-transform}

\subsection{Counters and transformed values}
\label{subsec:counters-transform}

Let $p_j$ be the $j$th prime and put $p_0=0$.  For real $x\ge1$, define
\begin{align}
N_+(x)&=\#\{n\le x:\mu(n)=1\},
\label{eq:Nplus-definition}\\
N_-(x)&=\#\{n\le x:\mu(n)=-1\},
\label{eq:Nminus-definition}\\
Q(x)&=N_+(x)+N_-(x),
\label{eq:Q-definition}\\
M(x)&=N_+(x)-N_-(x).
\label{eq:M-counter-identity}
\end{align}
The last definition agrees with \eqref{eq:mertens-definition}.  In
particular,
\begin{equation}
N_\pm(x)=\frac{Q(x)\pm M(x)}2.
\label{eq:counters-QM}
\end{equation}

\begin{definition}[Ordinal-prime transform]
\label{def:S-transform}
The sign-specific ordinal-prime transform and its summatory function are
\begin{align}
S(n)&=
\begin{cases}
p_{N_+(n)},&\mu(n)=1,\\
-p_{N_-(n)},&\mu(n)=-1,\\
1,&\mu(n)=0,
\end{cases}
\label{eq:S-definition}\\
SS(x)&=\sum_{n\le x}S(n).
\label{eq:SS-definition}
\end{align}
\end{definition}

Thus the first positive M\"obius value receives $p_1$, the second receives
$p_2$, and so on; the negative occurrences receive the negatives of the same
prime list on their own clock.  The value $1$ at nonsquarefree integers is part
of the definition and must be retained.

Put
\begin{equation}
\calP(m)=\sum_{j=1}^{m}p_j,
\qquad \calP(0)=0.
\label{eq:prime-prefix-definition}
\end{equation}

\begin{proposition}[Closed form for $SS$]
\label{prop:SS-closed-form}
For every real $x\ge1$,
\begin{equation}
SS(x)=\lfloor x\rfloor-Q(x)
+\calP\bigl(N_+(x)\bigr)-\calP\bigl(N_-(x)\bigr).
\label{eq:SS-closed-form}
\end{equation}
Consequently, with
\begin{equation}
T(x)=SS(x)-\lfloor x\rfloor+Q(x),
\label{eq:T-definition}
\end{equation}
one has
\begin{equation}
T(x)=\calP\bigl(N_+(x)\bigr)-\calP\bigl(N_-(x)\bigr).
\label{eq:T-prime-block}
\end{equation}
\end{proposition}

\begin{proof}
There are $\lfloor x\rfloor-Q(x)$ nonsquarefree integers up to $x$, each
contributing $1$.  The positive squarefree occurrences contribute the prime
prefix $p_1+\cdots+p_{N_+(x)}$, and the negative occurrences contribute the
negative prime prefix $-(p_1+\cdots+p_{N_-(x)})$.  Their sum is
\eqref{eq:SS-closed-form}; subtracting the zero row gives
\eqref{eq:T-prime-block}.
\end{proof}

\subsection{Ordinal and prime-displacement parts}
\label{subsec:ordinal-displacement}

For a squarefree $n$, let $j(n)$ denote its occurrence number inside its own
sign class.  Define
\begin{equation}
w_j=p_j-j,\qquad
D_w(m)=\sum_{j=1}^{m}w_j,\qquad D_w(0)=0.
\label{eq:w-Dw-definition}
\end{equation}

\begin{lemma}[Exact ordinal coefficient sum]
\label{lem:ordinal-coefficient-sum}
For every real $x\ge1$,
\begin{equation}
\sum_{\substack{n\le x\\ \mu(n)\ne0}}
\mu(n)j(n)=\frac{M(x)(Q(x)+1)}2.
\label{eq:ordinal-coefficient-sum}
\end{equation}
\end{lemma}

\begin{proof}
The positive occurrences contribute $1+\cdots+N_+$ and the negative
occurrences contribute $-(1+\cdots+N_-)$.  Therefore
\begin{align}
\sum_{\substack{n\le x\\ \mu(n)\ne0}}\mu(n)j(n)
&=\frac{N_+(N_++1)-N_-(N_-+1)}2
\label{eq:ordinal-sum-first}\\
&=\frac{(N_+-N_-)(N_++N_-+1)}2
\label{eq:ordinal-sum-second}\\
&=\frac{M(x)(Q(x)+1)}2.
\label{eq:ordinal-sum-third}
\end{align}
\end{proof}

\begin{definition}[Prime-displacement partial sum]
\label{def:W-partial-sum}
Set
\begin{equation}
W(x)=\sum_{\substack{n\le x\\ \mu(n)\ne0}}
\mu(n)w_{j(n)}.
\label{eq:W-definition}
\end{equation}
\end{definition}

\begin{proposition}[Exact recovery of $W$ from $SS$]
\label{prop:W-from-SS}
For every real $x\ge1$,
\begin{align}
W(x)&=D_w\bigl(N_+(x)\bigr)-D_w\bigl(N_-(x)\bigr)
\label{eq:W-Dw}\\
&=SS(x)-\lfloor x\rfloor+Q(x)
-\frac{M(x)(Q(x)+1)}2.
\label{eq:W-SS}
\end{align}
\end{proposition}

\begin{proof}
The first equality follows by summing $w_j$ on the two occurrence clocks.
Since $p_j=j+w_j$, the signed-prime contribution $T(x)$ is the sum of the
ordinal contribution in Lemma~\ref{lem:ordinal-coefficient-sum} and $W(x)$.
Use \eqref{eq:T-definition} and rearrange.
\end{proof}

\begin{proposition}[Asymptotic amplification]
\label{prop:W-asymptotic}
Let $c_0=3/\pi^2$.  As $x\to\infty$,
\begin{equation}
W(x)=c_0x\log x\,M(x)\bigl(1+o(1)\bigr),
\label{eq:W-asymptotic}
\end{equation}
with equality of both sides when $M(x)=0$ understood before taking the
relative error.
\end{proposition}

\begin{proof}
The prime number theorem gives $p_j\sim j\log j$ and $M(x)=o(x)$, while
elementary squarefree counting gives
\begin{equation}
Q(x)=\frac6{\pi^2}x+O(\sqrt{x})=2c_0x+O(\sqrt{x}).
\label{eq:Q-asymptotic}
\end{equation}
Hence both $N_+(x)$ and $N_-(x)$ are $c_0x(1+o(1))$.  If $M(x)>0$,
\eqref{eq:W-Dw} is the block sum of $w_j$ over
$N_-(x)<j\le N_+(x)$; if $M(x)<0$, it is the negative of the reversed
block.  Every index in that block is $c_0x(1+o(1))$, uniformly, and
\begin{equation}
w_j=p_j-j=c_0x\log x\bigl(1+o(1)\bigr)
\label{eq:w-uniform-block}
\end{equation}
uniformly on the block.  There are $|M(x)|$ terms and their common orientation
is $\sgn M(x)$, proving \eqref{eq:W-asymptotic}.
\end{proof}

\section{The Dirichlet series \texorpdfstring{$\alpha$}{alpha} and its exact decomposition}
\label{sec:alpha-decomposition}

\begin{definition}[The series $\alpha$ and $\delta$]
\label{def:alpha-delta}
Define
\begin{align}
\alpha(s)&=\sum_{n=1}^{\infty}\frac{S(n)}{n^s},
\label{eq:alpha-definition}\\
\delta(s)&=
\sum_{\substack{n\ge1\\ \mu(n)\ne0}}
\frac{\mu(n)w_{j(n)}}{n^s}.
\label{eq:delta-definition}
\end{align}
The analytic function $\alpha(s)$ is not the ordinal position $\alpha_j$
introduced later.
\end{definition}

Since $p_j\ll j\log(2j)$ and $j(n)\le n$, both series in
Definition~\ref{def:alpha-delta} converge absolutely in $\Re s>2$.  This is
the unconditional half-plane in which termwise regrouping will first be used.

\begin{proposition}[Four-block decomposition]
\label{prop:alpha-four-blocks}
For $\Re s>2$,
\begin{equation}
\alpha(s)=A(s)+B(s)+C(s)+\delta(s),
\label{eq:alpha-ABCD}
\end{equation}
where
\begin{align}
A(s)&=\zeta(s)-\frac{\zeta(s)}{\zeta(2s)},
\label{eq:A-definition}\\
B(s)&=c_0\frac{s}{(s-1)\zeta(s-1)},
\label{eq:B-definition}\\
C(s)&=\frac{s}{2}\int_1^\infty
M(x)\bigl(Q(x)+1-2c_0x\bigr)x^{-s-1}\,dx.
\label{eq:C-definition}
\end{align}
The functions $A$ and $B$ are meromorphic on $\C$.  The integral defining
$C$ converges locally uniformly and defines a holomorphic function in
$\Re s>3/2$.
\end{proposition}

\begin{proof}
The nonsquarefree row of $S$ equals
\begin{align}
\sum_{\mu(n)=0}\frac1{n^s}
&=\sum_{n\ge1}\frac1{n^s}
-\sum_{n\ge1}\frac{\mu(n)^2}{n^s}
\label{eq:zero-row-first}\\
&=\zeta(s)-\frac{\zeta(s)}{\zeta(2s)}=A(s),
\label{eq:zero-row-second}
\end{align}
where the second equality uses the Euler product for the squarefree indicator.
All series are absolutely convergent for $\Re s>2$, in fact for this step
$\Re s>1$ suffices.

For squarefree $n$, the identity $p_{j(n)}=j(n)+w_{j(n)}$ gives
\begin{equation}
S(n)=\mu(n)j(n)+\mu(n)w_{j(n)}.
\label{eq:S-ordinal-residual}
\end{equation}
The second term produces $\delta$.  Let
\begin{equation}
J(x)=\sum_{\substack{n\le x\\ \mu(n)\ne0}}\mu(n)j(n)
=\frac{M(x)(Q(x)+1)}2.
\label{eq:J-definition}
\end{equation}
Finite Abel summation followed by passage to infinity yields
\begin{equation}
\sum_{\substack{n\ge1\\ \mu(n)\ne0}}
\frac{\mu(n)j(n)}{n^s}
=s\int_1^\infty J(x)x^{-s-1}\,dx,
\label{eq:J-Mellin}
\end{equation}
because $J(x)=O(x^2)$ and $J(X)X^{-s}\to0$ for $\Re s>2$.
Split $Q(x)+1$ as $2c_0x+(Q(x)+1-2c_0x)$.  The remainder is exactly
$C(s)$.  For the main term, put $z=s-1$ and use
\begin{equation}
\int_1^\infty M(x)x^{-z-1}\,dx
=\frac1{z\zeta(z)},\qquad \Re z>1,
\label{eq:M-mellin}
\end{equation}
which follows by Abel summation from \eqref{eq:mobius-zeta}.  The main term is
therefore $B(s)$.

Finally, \eqref{eq:Q-asymptotic} and the unconditional bound $|M(x)|\le x$
give
\begin{equation}
\left|M(x)\bigl(Q(x)+1-2c_0x\bigr)x^{-s-1}\right|
\ll x^{1/2-\Re s}.
\label{eq:C-majorant}
\end{equation}
The majorant is integrable precisely when $\Re s>3/2$.  On every compact
subset of that half-plane it is uniform, so dominated convergence for
holomorphic parameter integrals proves holomorphy of $C$ there.
\end{proof}

\subsection{Singularities and the noncancellation warning}
\label{subsec:singularities}

\begin{proposition}[Singularities of the explicit shifted block]
\label{prop:B-singularities}
If $\rho$ is a nontrivial zero of $\zeta$ of multiplicity $m$, then $B$ has
a pole of order $m$ at $s=1+\rho$.  For a simple zero,
\begin{equation}
\Res_{s=1+\rho}B(s)
=c_0\frac{1+\rho}{\rho\zeta'(\rho)}.
\label{eq:B-residue}
\end{equation}
The function $B$ also has a simple pole at $s=1$ with residue $-2c_0$, poles
at $s=1-2k$ arising from the trivial zeros, and a zero rather than a pole at
$s=2$.
\end{proposition}

\begin{proof}
At a zero $\rho$, factor $\zeta(s-1)$ locally and note that
$c_0s/(s-1)$ is finite and nonzero at $1+\rho$.  Formula
\eqref{eq:B-residue} follows from the linear factor at a simple zero.  At
$s=1$, use $\zeta(0)=-1/2$.  At $s=2$, the pole of $\zeta(s-1)$ makes its
reciprocal vanish.  The trivial zeros of $\zeta(s-1)$ occur at $s=1-2k$.
\end{proof}

\begin{warning}[No automatic poles for the full $\alpha$]
\label{warn:alpha-poles}
Proposition~\ref{prop:B-singularities} locates poles of $B$, not of the full
$\alpha$.  Identity \eqref{eq:alpha-ABCD} is initially established only in
$\Re s>2$.  A pole of $B$ cannot be assigned to $\alpha$ outside that common
domain without first continuing the other terms and proving noncancellation.
In particular, neither $\delta$ nor $C$ may be silently absorbed into an
unspecified holomorphic remainder.
\end{warning}

\section{Squarefree time, ordinal pairs, and finite boundaries}
\label{sec:squarefree-pairs}

Let $s_1<s_2<\cdots$ enumerate the squarefree positive integers, and define
\begin{equation}
\varepsilon_q=\mu(s_q),\qquad
R(q)=\sum_{r=1}^{q}\varepsilon_r=M(s_q),\qquad R(0)=0.
\label{eq:squarefree-clock}
\end{equation}
The occurrence counters at squarefree time are
\begin{equation}
u_q=\frac{q+R(q)}2,qquad
v_q=\frac{q-R(q)}2.
\label{eq:uv-counters}
\end{equation}

Let $\alpha_j$ and $\beta_j$ denote the squarefree-time positions of the
$j$th positive and negative signs, respectively, and put
\begin{equation}
a_j=s_{\alpha_j},\qquad b_j=s_{\beta_j},
\qquad d_j=\beta_j-\alpha_j.
\label{eq:ordinal-pairs}
\end{equation}

\begin{proposition}[Pair representation]
\label{prop:pair-representation}
For $\Re s>2$,
\begin{equation}
\delta(s)=\sum_{j=1}^{\infty}w_j
\left(a_j^{-s}-b_j^{-s}\right).
\label{eq:delta-pair-representation}
\end{equation}
Moreover, for every complex $s$ and every $j$,
\begin{equation}
a_j^{-s}-b_j^{-s}
=s\int_{a_j}^{b_j}x^{-s-1}\,dx,
\label{eq:oriented-pair-integral}
\end{equation}
where the orientation of the integral is retained.
\end{proposition}

\begin{proof}
Absolute convergence in $\Re s>2$ permits matching the $j$th positive and
negative occurrences.  Their residual coefficients are $w_j$ and $-w_j$,
which gives \eqref{eq:delta-pair-representation}.  The derivative of
$x^{-s}$ is $-sx^{-s-1}$, proving \eqref{eq:oriented-pair-integral}; the
formula also holds at $s=0$ by direct evaluation.
\end{proof}

Define the coefficient at squarefree time and its cumulative sum by
\begin{align}
c_q&=\varepsilon_q w_{j(s_q)},
\label{eq:cq-definition}\\
U(q)&=\sum_{r=1}^{q}c_r
=D_w(u_q)-D_w(v_q).
\label{eq:U-definition}
\end{align}

\begin{proposition}[Open-pair block and $SS$ readout]
\label{prop:U-open-block}
For every $q\ge1$,
\begin{align}
U(q)&=\sgn R(q)
\sum_{j=\min(u_q,v_q)+1}^{\max(u_q,v_q)}w_j,
\label{eq:U-open-block}\\
U(q)&=SS(s_q)-s_q+q-\frac{q+1}{2}R(q).
\label{eq:U-from-SS}
\end{align}
An empty sum in \eqref{eq:U-open-block} is zero.
\end{proposition}

\begin{proof}
The difference $D_w(u_q)-D_w(v_q)$ is the oriented block in
\eqref{eq:U-open-block}.  It consists exactly of ordinal pairs with one
endpoint at or before cut $q$ and the other endpoint after it.  Formula
\eqref{eq:U-from-SS} is \eqref{eq:W-SS} evaluated at $x=s_q$, using
$Q(s_q)=q$ and $M(s_q)=R(q)$.
\end{proof}

\begin{lemma}[Finite Abel identity with the endpoint term]
\label{lem:finite-Abel}
For $Q\ge1$ and every $s\in\C$, let
\begin{equation}
\delta_Q(s)=\sum_{q=1}^{Q}c_qs_q^{-s}.
\label{eq:deltaQ-definition}
\end{equation}
Then
\begin{equation}
\delta_Q(s)=U(Q)s_Q^{-s}
+\sum_{q=1}^{Q-1}U(q)
\left(s_q^{-s}-s_{q+1}^{-s}\right).
\label{eq:finite-Abel-identity}
\end{equation}
\end{lemma}

\begin{proof}
Since $c_q=U(q)-U(q-1)$, with $U(0)=0$, finite summation gives
\begin{align}
\sum_{q=1}^{Q}c_qs_q^{-s}
&=\sum_{q=1}^{Q}\bigl(U(q)-U(q-1)\bigr)s_q^{-s}
\label{eq:finite-Abel-first}\\
&=U(Q)s_Q^{-s}
+\sum_{q=1}^{Q-1}U(q)
\bigl(s_q^{-s}-s_{q+1}^{-s}\bigr).
\label{eq:finite-Abel-second}
\end{align}
No infinite interchange occurs in this proof.
\end{proof}

\begin{remark}[Meaning of the boundary]
\label{rem:unfinished-boundary}
The term $U(Q)s_Q^{-s}$ is the contribution of all ordinal pairs still open at
the final cut.  Removing it replaces the original order by a completed-pair
grouping.  Such a replacement is not legitimate in the RH transfer unless the
boundary is independently shown to vanish.
\end{remark}

For $\Re s>2$, the bound $|U(q)|\ll s_q^2\log(2s_q)$ makes the infinite
cut representation absolutely convergent, and Lemma~\ref{lem:finite-Abel}
then yields
\begin{equation}
\delta(s)=\sum_{q=1}^{\infty}U(q)
\left(s_q^{-s}-s_{q+1}^{-s}\right).
\label{eq:infinite-cut-representation}
\end{equation}
This passage is not used to define the finite CMZ observable; the finite
identity \eqref{eq:finite-Abel-identity} is the primary formula.

\section{The CMZ multiscale mechanism: a proof-architecture summary}
\label{sec:CMZ-summary}

\subsection{Euler vorticity and hyperbolic stretching}
\label{subsec:Euler-stretching}

The forced incompressible Euler equations and their vorticity formulation are
\begin{align}
\partial_tu+(u\cdot\nabla)u+\nabla p&=f,
\qquad \nabla\cdot u=0,
\label{eq:forced-Euler}\\
\partial_t\omega+(u\cdot\nabla)\omega
&=(\omega\cdot\nabla)u+\nabla\times f,
\qquad \omega=\nabla\times u.
\label{eq:forced-vorticity}
\end{align}
The CMZ constructions arrange an older vorticity field whose velocity is
approximately linear near the origin.  A younger, much smaller oscillatory
packet is transported and stretched by this hyperbolic linearisation
\cite{CordobaMartinezZoroa2023,CordobaMartinezZoroaZheng2023}.

In the forced non-axisymmetric construction, a parent scale $N$ and child
scale $L$ are related by
\begin{equation}
L^{1/2}=exp\left(
\int_{1-N^{-\eps/2}}^1
\partial_1u_{N,1}(0,t)\,dt\right),
\label{eq:CMZ-scale-relation}
\end{equation}
under a positive strain bound of size $N^\eps$.  Writing
\begin{equation}
K_N(t)=\exp\left(\int_t^1\partial_1u_{N,1}(0,s)\,ds\right),
\qquad
b(t)=\frac{\log K_N(t)}{\log L},
\label{eq:CMZ-K-b}
\end{equation}
the transported child packet has anisotropic widths of orders
\begin{equation}
L^{-1/2-b+\eps},\qquad
L^{-1/2-\eps+b},\qquad
L^{-1/2},
\label{eq:CMZ-support-widths}
\end{equation}
and oscillation frequency of order $L$.

\subsection{Parametrix, cancellation, and gluing}
\label{subsec:CMZ-proof-stages}

The proof architecture in \cite{CordobaMartinezZoroa2023} has the following
linked stages.
\begin{enumerate}[label=\textup{C\arabic*.},leftmargin=*]
\item A local Biot--Savart parametrix approximates the two leading velocity
components generated by the packet.  For spatial derivatives of order
$r=0,1$, the quoted error has the form
\begin{equation}
\|\widetilde u_j-u_j\|_{C^r}
\le C_\kappa L^{2\eps+r-3/2+\kappa}.
\label{eq:CMZ-parametrix-error}
\end{equation}
\item The leading pieces of
$\widetilde u_2\partial_2\omega_1$ and
$\widetilde u_3\partial_3\omega_1$ cancel.  The proof printed after the
stated lemma safely yields
\begin{equation}
\|u(\omega)\cdot\nabla\omega\|_{C^\vartheta}
+\|\omega\cdot\nabla u(\omega)\|_{C^\vartheta}
\lesssim L^{4\eps-1/2+\vartheta}.
\label{eq:CMZ-safe-quadratic}
\end{equation}
The exponent $4\eps$ in \eqref{eq:CMZ-safe-quadratic}, rather than the
$3\eps$ appearing in the lemma statement, is the exponent obtained after the
paper chooses its auxiliary loss parameter equal to $\eps$.
\item Oscillation and spatial separation make the child's velocity decay
faster than any prescribed power outside its anisotropic support.  Odd parity
and Taylor expansion control interactions with the parent field.
\item Replacing the parent velocity by its local linear model produces a
controlled forcing error.  With
$\vartheta=1/2-6\eps$, the safe power in
\eqref{eq:CMZ-safe-quadratic} becomes $L^{-2\eps}$.
\item The induction alternates the expanding direction and adds a child
whose strain at the origin is comparable to $L^\eps$, while its forcing
increment is bounded by
\begin{equation}
\|F_{\mathrm{child}}\|_{C^{1/2-6\eps}}
\lesssim L^{-2\eps}.
\label{eq:CMZ-force-increment}
\end{equation}
Choosing scales sufficiently separated makes the forcing increments summable.
\end{enumerate}

These steps summarise the actual proof mechanism; they do not assert that an
arbitrary prescribed time-dependent strain can be realised with constants
independent of its magnitude.

\subsection{The separated-layer amplitude system}
\label{subsec:CMZZ-ODE}

In the unforced separated-layer construction of
\cite{CordobaMartinezZoroaZheng2023}, the height $x_k(t)>0$ of the $k$th
packet is governed, up to controlled geometric errors, by
\begin{equation}
\dot x_k(t)=x_k(t)\sum_{j<k}a_j(t)x_j(t),
\label{eq:CMZZ-ODE}
\end{equation}
where $a_j(t)$ is close to $1$ and, importantly, does not depend on the child
index $k$.  In the exact model $a_j=1$, writing $y_k=\log x_k$ gives
\begin{equation}
\dot y_k(t)=\sum_{j<k}x_j(t).
\label{eq:CMZZ-log-law}
\end{equation}
For terminal data $x_k(0)=A^k$, $A>1$, the model admits a uniform bound on
the time integral of each amplitude over a fixed backward interval:
\begin{equation}
\int_{-a}^{0}x_k(t)\,dt\le a,
\qquad a=A(1-e^{-a}).
\label{eq:CMZZ-integral-bound}
\end{equation}
This positivity is structural.  Encoding $\mu$ by allowing signed amplitudes
would destroy the hypothesis behind \eqref{eq:CMZZ-integral-bound}.  The
two-channel construction in Section~\ref{sec:CMZ-deformation} avoids that
problem by keeping both channels positive.

\subsection{Later multiscale developments}
\label{subsec:later-multiscale}

The 2026 IPM manuscript \cite{AlpogeBuckmasterCoiculescu2026} transports
oscillatory packets through nested Fourier cutoffs and proves joint
space--time smoothness of the force on the torus.  The 2026 Boussinesq
manuscript \cite{AlpogeBuckmaster2026} uses an exact growing wave, compactly
supported corrections, lifetime derivative estimates, and summable local
corrections to obtain smooth forcing in both equations.  These works confirm
the flexibility of the multiscale programme.  Their theorems concern different
PDEs and do not imply an arithmetic cancellation estimate for $U(q)$.

\section{The positive two-channel CMZ deformation}
\label{sec:CMZ-deformation}

The construction in this section is an \textup{[ADAPT]} of the logarithmic
stretching law \eqref{eq:CMZZ-log-law}.  It is an exact finite-dimensional
dynamical encoding, not yet an Euler solution.

Fix a real $\sigma>0$ and a cutoff $Q\ge1$.  Define positive clock lengths
\begin{equation}
\ell_q^{(Q)}(\sigma)=
\begin{cases}
s_q^{-\sigma}-s_{q+1}^{-\sigma},&1\le q<Q,\\
s_Q^{-\sigma},&q=Q.
\end{cases}
\label{eq:CMZ-clock-lengths}
\end{equation}
Since $s_1=1$, these lengths satisfy
\begin{equation}
\sum_{q=1}^{Q}\ell_q^{(Q)}(\sigma)=1.
\label{eq:CMZ-clock-total}
\end{equation}
Let $0=\tau_0<\tau_1<\cdots<\tau_Q=1$ be defined by
\begin{equation}
\tau_q=\sum_{r=1}^{q}\ell_r^{(Q)}(\sigma).
\label{eq:CMZ-clock-partition}
\end{equation}

Write $U_+(q)=\max\{U(q),0\}$ and
$U_-(q)=\max\{-U(q),0\}$.  Define the nonnegative, piecewise constant
strains
\begin{equation}
g_{\pm,Q,\sigma}(t)=U_\pm(q),
\qquad \tau_{q-1}\le t<\tau_q.
\label{eq:positive-strains}
\end{equation}

\begin{definition}[Paired CMZ-type channels]
\label{def:paired-channels}
Let $X_{\pm,Q,\sigma}$ solve
\begin{equation}
\dot X_{\pm,Q,\sigma}(t)
=g_{\pm,Q,\sigma}(t)X_{\pm,Q,\sigma}(t),
\qquad X_{\pm,Q,\sigma}(0)=1.
\label{eq:paired-channel-ODE}
\end{equation}
Their relative terminal deformation is
\begin{equation}
\mathscr C_Q(\sigma)=
\frac{X_{+,Q,\sigma}(1)}{X_{-,Q,\sigma}(1)}.
\label{eq:relative-deformation}
\end{equation}
\end{definition}

Both amplitudes in \eqref{eq:paired-channel-ODE} remain positive.  The
arithmetic sign occurs only through their ratio.

\begin{proposition}[Exact arithmetic readout]
\label{prop:exact-CMZ-readout}
For every $Q\ge1$ and every real $\sigma>0$,
\begin{equation}
\log\mathscr C_Q(\sigma)=\delta_Q(\sigma).
\label{eq:CMZ-readout}
\end{equation}
\end{proposition}

\begin{proof}
Solving the scalar equations gives
\begin{equation}
\log X_{\pm,Q,\sigma}(1)
=\sum_{q=1}^{Q}U_\pm(q)\ell_q^{(Q)}(\sigma).
\label{eq:channel-solution}
\end{equation}
Subtracting the two identities in \eqref{eq:channel-solution} yields
\begin{equation}
\log\mathscr C_Q(\sigma)
=\sum_{q=1}^{Q}U(q)\ell_q^{(Q)}(\sigma).
\label{eq:channel-difference}
\end{equation}
Insert \eqref{eq:CMZ-clock-lengths}.  The result is precisely the finite Abel
identity \eqref{eq:finite-Abel-identity}.
\end{proof}

\begin{remark}[Why real $\sigma$ is used]
\label{rem:real-sigma}
For nonreal $s$, the differences $s_q^{-s}-s_{q+1}^{-s}$ are complex and
cannot serve as positive time lengths.  Proposition
\ref{prop:exact-CMZ-readout} therefore uses real $\sigma$.  This loses no
information in the main equivalence, because convergence of $\delta(\sigma)$
for every real $\sigma>3/2$ is already equivalent to RH.  Complex continuation
is treated separately in Section~\ref{sec:analytic-criterion}.
\end{remark}

\section{The CMZ deformation equivalence with RH}
\label{sec:main-equivalence}

We now state the exact bridge.  It is a theorem about the explicitly defined
deformation in Definition~\ref{def:paired-channels}; it does not assume a PDE
realisation.

\begin{theorem}[Ordinal-prime CMZ deformation criterion]
\label{thm:CMZ-RH-equivalence}
The following statements are equivalent.
\begin{enumerate}[label=\textup{(\roman*)},leftmargin=*]
\item The Riemann hypothesis holds.
\item For every real $\sigma>3/2$, the sequence
$\delta_Q(\sigma)$ converges as $Q\to\infty$.
\item For every real $\sigma>3/2$, the positive relative deformation
$\mathscr C_Q(\sigma)$ converges to a finite nonzero limit.
\item For every real $\sigma>3/2$,
\begin{equation}
\lim_{X\to\infty}\ 
\sup_{Y\ge X}
\left|
\log\frac{\mathscr C_Y(\sigma)}{\mathscr C_X(\sigma)}
\right|=0.
\label{eq:CMZ-tail-criterion}
\end{equation}
\end{enumerate}
\end{theorem}

\begin{proof}
Because the coefficients $c_n=\mathbf1_{\mu(n)\ne0}\mu(n)w_{j(n)}$
vanish at every nonsquarefree integer, the squarefree cutoff
$\delta_Q(s)=\sum_{q\le Q}c_{s_q}s_q^{-s}$ is exactly the ordinary partial
sum through the integer $s_Q$.  Thus convergence on the squarefree clock is
equivalent to convergence in the original integer ordering; no rearrangement
is being made.

Assume first RH.  By the Littlewood criterion
\eqref{eq:littlewood-criterion}, for every $\eta>0$,
\begin{equation}
M(x)=O_\eta\left(x^{1/2+\eta}\right).
\label{eq:RH-M-bound}
\end{equation}
The block formula \eqref{eq:W-Dw}, the prime number theorem, and
\eqref{eq:RH-M-bound} imply
\begin{equation}
W(x)=O_\eta\left(x^{3/2+\eta}\log(2x)\right).
\label{eq:RH-W-bound}
\end{equation}
Fix $\sigma>3/2$ and choose $\eta$ with
$0<\eta<\sigma-3/2$.  Abel summation between integer cutoffs $X<Y$ gives
\begin{align}
\sum_{X<n\le Y}
\frac{\mathbf1_{\mu(n)\ne0}\mu(n)w_{j(n)}}{n^\sigma}
&=W(Y)Y^{-\sigma}-W(X)X^{-\sigma}
\label{eq:delta-tail-Abel-one}\\
&\quad+\sigma\int_X^YW(x)x^{-\sigma-1}\,dx.
\label{eq:delta-tail-Abel-two}
\end{align}
By \eqref{eq:RH-W-bound}, the endpoint terms tend to zero and the integral
tail is dominated by a constant times
\begin{equation}
\int_X^\infty x^{1/2+\eta-\sigma}\log(2x)\,dx,
\label{eq:delta-tail-majorant}
\end{equation}
which tends to zero because $1/2+\eta-\sigma<-1$.  Thus the original-order
series $\delta(\sigma)$ converges, proving (i)$\Rightarrow$(ii).

Conversely, suppose (ii).  Fix any $\epsilon>0$ and take
\begin{equation}
\sigma=\frac32+\frac{\epsilon}{2}.
\label{eq:sigma-choice}
\end{equation}
Convergence of $\sum c_nn^{-\sigma}$, where
$c_n=\mathbf1_{\mu(n)\ne0}\mu(n)w_{j(n)}$, permits Kronecker's lemma and
gives
\begin{equation}
W(x)=o\left(x^\sigma\right).
\label{eq:Kronecker-W}
\end{equation}
Combine \eqref{eq:Kronecker-W} with the two-sided asymptotic
\eqref{eq:W-asymptotic}.  At points where $M(x)\ne0$ this gives
\begin{equation}
M(x)=o\left(\frac{x^{\sigma-1}}{\log x}\right)
=o\left(\frac{x^{1/2+\epsilon/2}}{\log x}\right).
\label{eq:M-from-delta}
\end{equation}
At points where $M(x)=0$, the desired bound is automatic.  Therefore
$M(x)=O_\epsilon(x^{1/2+\epsilon})$.  Since $\epsilon>0$ was arbitrary,
\eqref{eq:littlewood-criterion} gives RH.  This proves
(ii)$\Rightarrow$(i).

Proposition~\ref{prop:exact-CMZ-readout} gives
$\mathscr C_Q(\sigma)=\exp(\delta_Q(\sigma))$.  Since the real exponential
is a homeomorphism from $\R$ to $(0,\infty)$, (ii) is equivalent to (iii).
Finally,
\begin{equation}
\log\frac{\mathscr C_Y(\sigma)}{\mathscr C_X(\sigma)}
=\delta_Y(\sigma)-\delta_X(\sigma),
\label{eq:tail-log-delta}
\end{equation}
so the Cauchy criterion proves the equivalence of (ii) and (iv).
\end{proof}

\begin{remark}[The first RH-equivalent node]
\label{rem:first-equivalent-node}
Definitions of $S$, $SS$, $U$, and the finite deformation are unconditional.
The first statement in the dependency chain that is equivalent to RH is the
tail condition \eqref{eq:CMZ-tail-criterion}.  Consequently, an Euler estimate
that implies this condition with an unattenuated arithmetic normalisation is
an RH-strength estimate; it cannot be counted as a routine consequence of
the existence of a CMZ gluing.
\end{remark}

\section{The analytic \texorpdfstring{$\alpha$}{alpha} criterion and singularity audit}
\label{sec:analytic-criterion}

Let
\begin{equation}
\mathcal H=\{s\in\C:\Re s>3/2\}.
\label{eq:H-half-plane}
\end{equation}

\begin{theorem}[Shifted reciprocal-zeta criterion]
\label{thm:analytic-alpha-criterion}
The following are equivalent.
\begin{enumerate}[label=\textup{(\roman*)},leftmargin=*]
\item RH holds.
\item $B$ is holomorphic on $\mathcal H$.
\item The function $\alpha-\delta$, initially defined on $\Re s>2$, admits a
holomorphic continuation to $\mathcal H$.
\end{enumerate}
\end{theorem}

\begin{proof}
By \eqref{eq:B-definition}, poles of $B$ in $\mathcal H$ are exactly the
points $1+\rho$ arising from zeros $\rho$ of $\zeta$ with
$\Re\rho>1/2$; the elementary poles listed in Proposition
\ref{prop:B-singularities} lie outside $\mathcal H$.  Thus (i) is equivalent
to (ii), using the functional-equation symmetry of nontrivial zeros.

In $\mathcal H$, the function $C$ is holomorphic by Proposition
\ref{prop:alpha-four-blocks}.  Also $A$ is holomorphic there: its possible
denominator singularities arise from zeros of $\zeta(2s)$, which have
$\Re s<1/2$, while the pole at $s=1$ is outside $\mathcal H$.  On
$\Re s>2$, identity \eqref{eq:alpha-ABCD} gives
\begin{equation}
\alpha(s)-\delta(s)=A(s)+B(s)+C(s).
\label{eq:alpha-minus-delta}
\end{equation}
Therefore the right-hand side extends holomorphically to $\mathcal H$ exactly
when $B$ does.  Uniqueness follows from the identity theorem, proving the
equivalence of (ii) and (iii).
\end{proof}

\begin{warning}[Two distinct routes]
\label{warn:two-routes}
Theorem~\ref{thm:CMZ-RH-equivalence} concerns original-order convergence of
$\delta_Q(\sigma)$ on the real axis.  Theorem
\ref{thm:analytic-alpha-criterion} concerns continuation of
$\alpha-\delta$.  A regrouped continuation of $\delta$ does not by itself
prove original-order convergence, and holomorphy of $\delta$ alone does not
cancel a pole of $B$ in the full $\alpha$.
\end{warning}

\begin{corollary}[Abscissa statement]
\label{cor:abscissa}
Let
\begin{equation}
\Theta=\sup_{\zeta(\rho)=0,\ 0<\Re\rho<1}\Re\rho.
\label{eq:Theta-definition}
\end{equation}
For the original ordering of the Dirichlet series,
\begin{equation}
\sigma_c(\delta)=\sigma_c(\alpha)=1+\Theta.
\label{eq:abscissae}
\end{equation}
\end{corollary}

\begin{proof}
The classical growth-abscissa relation for $M$ gives growth exponent
$\Theta$ \cite{Titchmarsh1986}.  Proposition~\ref{prop:W-asymptotic} raises
that exponent by one for $W$, so the abscissa of $\delta$ is $1+\Theta$.
The signed-prime part of $SS$ has the same exponent, whereas the zero row has
partial sums $O(x)$; hence $\alpha$ has the same abscissa.  This argument is
about ordinary ordering and does not assert a natural boundary or a line of
poles.
\end{proof}

\section{Conditional realisation by genuine CMZ towers}
\label{sec:PDE-realisation}

The exact criterion above uses the CMZ logarithmic law, but its prescribed
piecewise strain $g_{\pm,Q,\sigma}$ is not automatically generated by a
solution of the Euler equation.  We isolate the additional statement needed
for a genuine PDE realisation.

\begin{definition}[Admissible CMZ readout]
\label{def:admissible-readout}
For fixed real $\sigma>3/2$, an admissible CMZ readout is a family of positive
PDE deformation ratios $\widetilde{\mathscr C}_Q(\sigma)$ for which
\begin{equation}
\log\widetilde{\mathscr C}_Q(\sigma)
=\delta_Q(\sigma)+r_Q(\sigma),
\label{eq:PDE-readout-error}
\end{equation}
where the realisation error satisfies
\begin{equation}
\lim_{X\to\infty}\sup_{Y\ge X}
|r_Y(\sigma)-r_X(\sigma)|=0.
\label{eq:readout-error-Cauchy}
\end{equation}
The coefficient of $\delta_Q$ in \eqref{eq:PDE-readout-error} is exactly one
and is independent of $Q$ and of all spatial carrier frequencies.
\end{definition}

\begin{theorem}[Conditional Euler-CMZ transfer]
\label{thm:conditional-PDE-transfer}
Assume that an admissible CMZ readout exists for every real
$\sigma>3/2$.  Then RH is equivalent to
\begin{equation}
\lim_{X\to\infty}\sup_{Y\ge X}
\left|
\log\frac{\widetilde{\mathscr C}_Y(\sigma)}
{\widetilde{\mathscr C}_X(\sigma)}
\right|=0
\quad\text{for every real }\sigma>3/2.
\label{eq:PDE-CMZ-tail}
\end{equation}
\end{theorem}

\begin{proof}
Subtract \eqref{eq:PDE-readout-error} at $Y$ and $X$ to obtain
\begin{align}
\log\frac{\widetilde{\mathscr C}_Y(\sigma)}
{\widetilde{\mathscr C}_X(\sigma)}
&=\delta_Y(\sigma)-\delta_X(\sigma)
\label{eq:PDE-tail-decomposition-one}\\
&\quad+r_Y(\sigma)-r_X(\sigma).
\label{eq:PDE-tail-decomposition-two}
\end{align}
Condition \eqref{eq:readout-error-Cauchy} shows that
\eqref{eq:PDE-CMZ-tail} holds exactly when the $\delta$ tails satisfy the
Cauchy condition.  Apply Theorem~\ref{thm:CMZ-RH-equivalence}.
\end{proof}

\begin{question}[Paired-reunion problem]
\label{q:paired-reunion}
Can one construct two genuine CMZ towers from the ordinal pairs
$+p_j,-p_j$ so that their principal Biot--Savart strains cancel after both
members have appeared, while the transported reunion defects satisfy
\eqref{eq:readout-error-Cauchy} with the exact, frequency-independent
normalisation in \eqref{eq:PDE-readout-error}?
\end{question}

Question~\ref{q:paired-reunion} is \textup{[OPEN]}.  Fixed-duration pairs are
compatible with local transport estimates, but long pairs can undergo large
hyperbolic deformation before reunion.  Allowing the carrier scale to absorb
$|U(q)|$ does not answer the question: it changes the readout coefficient and
can make any finite datum appear small.

\subsection{Three attempted insertions of the ordinal-prime data}
\label{subsec:three-attempts}

The two-channel model of Section~\ref{sec:CMZ-deformation} reads $U(q)$ off
an already-defined scalar ODE; it does not feed $SS$ into a genuine solution
of a forced Euler or IPM equation.  Three routes toward such a feeding are
recorded here, because each rules out a naive fix and narrows what a correct
insertion must do.

\textbf{(a) Rescaling the amplitude of a single packet.}  In the forced
construction of \cite{CordobaMartinezZoroa2023} (its Theorem~5.1) and in its
space--time-smooth extension \cite{AlpogeBuckmasterCoiculescu2026} (its
Proposition~5.13, eq.~(101), and Corollary~4.10), the leading amplitude of a
child packet is not a free parameter.  In \cite{CordobaMartinezZoroa2023},
the parent hypothesis pins $\partial_{x_1}u_1(\bar\omega_1)$ to
$[C_0/4,C_0]N^\varepsilon$ over a window of length $N^{-\varepsilon/2}$,
forcing the next scale $M^{1/2}=\exp\big(\Theta(N^{\varepsilon/2})\big)$; in
\cite{AlpogeBuckmasterCoiculescu2026}, the amplitude $a_{n+1}(t)$ solves an
ODE driven by the parent field and is normalised at the terminal time,
$a_{n+1}(1)=1$, which forces the seed value $a_{n+1}(s_{n+1})$ to be
exponentially small.  Multiplying an already-built exact solution by an
external factor changes the force by a term of the same nonlinear order as
the self-interaction already accounted for in \eqref{eq:CMZ-force-increment},
and no estimate in either source covers that modified force.
\textup{[OPEN, ruled out as stated]}: neither source offers a free per-event
amplitude dial.

\textbf{(b) Continuous exponent insertion.}  In
\cite{AlpogeBuckmasterCoiculescu2026}, the scale exponent
$\beta_{n+1}\in[\beta_n/2,\beta_n]$ of its Proposition~5.13 is a genuine
interval, and the threshold $\Lambda_{\mathrm{sep}}$ of its Proposition~4.2
is proved uniform over that whole interval.  It is tempting to set
$\beta_{n(q)+1}/\beta_{n(q)}=\tfrac34\exp(\kappa\, c_qs_q^{-\sigma})\in[\tfrac12,1]$
and read $\delta_Q(\sigma)$ off $\log\beta_{n(Q)+1}$.  Checked against the
same paper's global assembly rather than against the one-step proposition
alone, this fails for a structural reason: the recursive exponent choice
that proves global existence there (its eqs.~(172)--(174)) never uses an
intermediate value of $\beta_{n+1}$ --- at every level it sets $\beta_{n+1}$
to exactly $\beta_n$ or exactly $\beta_n/2$, and the argument proving
$\beta_n\downarrow0$ (its Lemma~6.3) uses this finitely-many-halvings
dichotomy directly.  A continuously varying exponent schedule is therefore
outside what the cited global theorem proves, not an unverified special case
of it.  \textup{[OPEN]}: encoding $\delta_Q$ in $\beta_{n+1}$ this way needs
a new global-iteration theorem for a continuously varying exponent, which is
not in the cited sources.  It is also, on its own terms, the wrong kind of
fix: the same scheme accepts \emph{any} sequence $x_q\to0$, so proving the
extension would encode $SS$ without constraining it.

\textbf{(c) Insertion into the parent gradient.}  Writing
$P_{n+1}=S_{\mu_{n+1}}\rho^{(n)}$ for the actual parent of the next packet
(eq.~(107) of \cite{AlpogeBuckmasterCoiculescu2026}), the real amplitude law
there is
\begin{equation}
\frac{\dot a_{n+1}}{a_{n+1}}
=C_0\cos\alpha_{n+1}\,e^\perp(\alpha_{n+1})\cdot\nabla P_{n+1}(0,t),
\qquad a_{n+1}(1)=1.
\label{eq:CMZ-amplitude-real}
\end{equation}
Asking that
\begin{equation}
\nabla P_{n(q)+1}(0,t_q)
=\nabla P^{\mathrm{cascade}}_{n(q)+1}(0,t_q)+\mathbf v\,U(q)+E_q
\label{eq:parent-insertion-target}
\end{equation}
for a fixed vector $\mathbf v$ and a summable error, so that $U(q)$ enters
\eqref{eq:CMZ-amplitude-real} through an output of the construction rather
than through a hand-chosen parameter, is structurally the right kind of
target: $\nabla P_{n+1}$ is forced by the sum of earlier packets, not free.
It is also compatible with the actual size of the data: $c_q=O(q\log q)$
grows, and the native CMZ regime is one of growing packet amplitudes, not
decaying ones, unlike attempts (a) and (b) above.  No construction of the
auxiliary packets in \eqref{eq:parent-insertion-target} is given here or in
the sources consulted; it is recorded below as
Question~\ref{q:parent-insertion} and is \textup{[OPEN]}.

\textbf{(d) The single-tower amplitude ratio is impossible, not merely
unproven.}  A natural sharpening of (c) proposes reading $U(q)$ directly off
a single tower's amplitude ratio,
\begin{equation}
\log\frac{a_{q+1}(1)}{a_{q+1}(s_{q+1})}
=c_qs_q^{-\sigma}+e_q.
\label{eq:single-tower-attempt}
\end{equation}
This is excluded outright, not merely unverified.

\begin{proposition}[Sign obstruction for the single-tower amplitude ratio]
\label{prop:sign-obstruction}
For every admissible family in the sense of
\cite{AlpogeBuckmasterCoiculescu2026} and every $q$ past the threshold of
its Lemma~4.9,
\begin{equation}
\log\frac{a_{q+1}(1)}{a_{q+1}(s_{q+1})}\ge k_g\,\lambda_q^{\beta_q/8}>0,
\qquad k_g=\frac1{96\pi}
\label{eq:amplitude-gain-lower-bound}
\end{equation}
(its eq.~(102)).  On the other hand $c_q=\varepsilon_qw_{\alpha_q}$, with
$\varepsilon_q=\mu(s_q)\in\{\pm1\}$ changing sign infinitely often and
$w_j=p_j-j>0$ for every $j\ge1$, so $c_qs_q^{-\sigma}$ changes sign
infinitely often, unconditionally.  Consequently
\eqref{eq:single-tower-attempt} cannot hold with $e_q\to0$: its left side is
always positive while its right side is negative for infinitely many $q$.
\end{proposition}

\begin{proof}
Immediate from \eqref{eq:amplitude-gain-lower-bound} and the sign pattern
of $\mu$.
\end{proof}

This does not merely rule out an unverified special case; it contradicts an
already-proved lemma of \cite{AlpogeBuckmasterCoiculescu2026}.  A two-tower
difference (the architecture of Section~\ref{sec:CMZ-deformation}, now built
from genuine packets on each side) evades the sign problem, since a
difference of two positive quantities has no fixed sign.  It does not evade
the size problem, treated in Section~\ref{subsec:magnitude-obstruction}.

\subsection{A generalisation of the layered amplitude bound}
\label{subsec:generalised-lemma}

The unforced layered construction of \cite{CordobaMartinezZoroaZheng2023}
states its self-interaction bound (its Lemma~1) only for geometric initial
data $x_n(0)=A^n$, with a single ratio $A>1$ fixed for the whole tower.  The
following removes that restriction.  It is proved here and is not present in
the cited source.

\begin{lemma}[Bounded-ratio generalisation]
\label{lem:generalised-layer-bound}
Let $(B_n)_{n\ge0}$ be positive reals with $B_0=1$ and
\begin{equation}
r_n:=B_n/B_{n-1}\in(0,\bar r]\qquad(n\ge1),\qquad\bar r>1.
\label{eq:ratio-range}
\end{equation}
Suppose $x_n(t)>0$ solves $\dot x_n(t)=x_n(t)\sum_{j<n}a_j(t)x_j(t)$ with
$a_j(t)\ge1$ and $x_n(0)=B_n$.  Let $\zeta>0$ solve $\zeta=\bar r(1-e^{-\zeta})$.
Then, for every $n\ge0$,
\begin{equation}
\int_{-\zeta}^0x_n(t)\,dt\le\zeta.
\label{eq:generalised-layer-bound}
\end{equation}
\end{lemma}

\begin{proof}
Put $y_n=\ln x_n$ and $Z_n(T)=\int_{-T}^0x_n(t)\,dt$.  Since $a_n(t)\ge1$,
$\dot y_{n+1}(t)\ge\dot y_n(t)+e^{y_n(t)}$; integrating from $t\le0$ to $0$
and exponentiating,
\begin{equation}
x_{n+1}(t)\le r_{n+1}\,x_n(t)\exp\Big(-\int_t^0x_n(s)\,ds\Big).
\label{eq:generalised-layer-step}
\end{equation}
Integrating \eqref{eq:generalised-layer-step} over $t\in[-T,0]$ and using
$\frac{d}{dt}\exp(-\int_t^0x_n(s)ds)=x_n(t)\exp(-\int_t^0x_n(s)ds)$ gives
\begin{equation}
Z_{n+1}(T)\le r_{n+1}\big(1-e^{-Z_n(T)}\big).
\label{eq:generalised-layer-recursion}
\end{equation}
Since $x_0\equiv B_0=1$, $Z_0(T)=T$; take $T=\zeta$, so $Z_0(\zeta)=\zeta$.
If $Z_n(\zeta)\le\zeta$, then \eqref{eq:generalised-layer-recursion},
$r_{n+1}\le\bar r$, and monotonicity of $x\mapsto1-e^{-x}$ give
$Z_{n+1}(\zeta)\le\bar r(1-e^{-\zeta})=\zeta$.  Induction closes.
\end{proof}

\begin{remark}[Scope of Lemma~\ref{lem:generalised-layer-bound}]
\label{rem:generalised-layer-scope}
This replaces the single global ratio $A$ by a bounded, otherwise arbitrary
sequence of ratios $r_n$, so in principle a different target can be assigned
at each level.  It does not by itself insert $SS$: the construction of
\cite{CordobaMartinezZoroaZheng2023} is a finite-time blow-up at $t=0$,
driven by $x_n(0)\to\infty$ along the tower, whereas the quantity needed for
Theorem~\ref{thm:CMZ-RH-equivalence} is convergence, as $Q\to\infty$, of a
Cauchy tail --- the opposite regime.  Reusing this lemma for the
ordinal-prime problem therefore requires either a genuinely different,
non-blow-up time parametrisation, or keeping only the self-interaction
estimate and discarding the blow-up end state; neither is done here.
\end{remark}

\subsection{A quadratic-summability bound, and the size obstruction to a two-tower comparison}
\label{subsec:magnitude-obstruction}

The following unconditional estimate isolates exactly how small the
arithmetic target already is, independently of any PDE construction.

\begin{proposition}[Quadratic summability of the ordinal-prime target]
\label{prop:quadratic-summability}
Let $d_q(\sigma):=c_qs_q^{-\sigma}$.  Then $|d_q(\sigma)|\ll_\sigma
q^{1-\sigma}\log(2q)$, and
\begin{equation}
\sum_{q=1}^\infty|d_q(\sigma)|^2<\infty\qquad\text{for every }\sigma>\frac32.
\label{eq:quadratic-summability}
\end{equation}
\end{proposition}

\begin{proof}
Since $\alpha_q\le q$ and $p_j\ll j\log(2j)$ (prime number theorem),
$|c_q|=w_{\alpha_q}=p_{\alpha_q}-\alpha_q\ll q\log(2q)$.  Since $s_q\asymp q$
(density of squarefree integers), $|d_q(\sigma)|\ll q^{1-\sigma}\log(2q)$,
and \eqref{eq:quadratic-summability} converges exactly when $2\sigma-2>1$,
i.e.\ $\sigma>3/2$.
\end{proof}

This uses only the prime number theorem and squarefree density --- no input
equivalent to RH --- and shows the target sequence is already small on its
own terms.  The next result shows this smallness works against, not for, a
two-tower realisation built from activated CMZ packets.

\begin{proposition}[Size obstruction for a two-tower comparison]
\label{prop:size-obstruction}
Let $\mathcal T^+,\mathcal T^-$ be two admissible families in the sense of
\cite{AlpogeBuckmasterCoiculescu2026}, each satisfying its global iteration
(so $\beta_q^\pm\downarrow0$ and $x_q^\pm:=\beta_q^\pm\log\lambda_q^\pm
\to\infty$, its eq.~(176)).  If a layer is activated in either tower at
step $q$, its amplitude-gain contribution to that tower's log-amplitude is,
by \eqref{eq:amplitude-gain-lower-bound}, at least
$k_g\exp(x_q^\pm/8)\to\infty$.  An insertion scheme in which exactly one
tower is activated at each step cannot satisfy
\begin{equation}
\left|
\log\frac{a^+_{q+1}(1)}{a^+_{q+1}(s_{q+1}^+)}
-\log\frac{a^-_{q+1}(1)}{a^-_{q+1}(s_{q+1}^-)}
\right|
=|d_q(\sigma)|+O(\eta_q),\qquad \eta_q\to0,
\label{eq:size-obstruction}
\end{equation}
for $q$ past the threshold where the inactive side contributes $0$ and the
active side's contribution exceeds any fixed bound, since the left side
then diverges while the right side, by
Proposition~\ref{prop:quadratic-summability}, tends to $0$.  If instead both
towers are activated at every step, \eqref{eq:size-obstruction} demands that
two independently-evolving, individually unbounded quantities agree to
within a vanishing relative error at every scale --- a synchronisation
condition supplied by no bound in
\cite{CordobaMartinezZoroa2023,CordobaMartinezZoroaZheng2023,AlpogeBuckmasterCoiculescu2026}.
\end{proposition}

\begin{remark}[What this rules out and what it does not]
\label{rem:size-obstruction-scope}
Proposition~\ref{prop:size-obstruction} rules out any two-tower realisation
built from the \emph{activated-packet} mechanism of
\cite{AlpogeBuckmasterCoiculescu2026} --- the crossing of its Lemma~4.8
followed by the gain of its Lemma~4.9 --- because that mechanism is, by
design, an amplifier: every activation is forced to contribute an
unboundedly growing amount, the opposite of the boundedly small, freely
signed contribution that Proposition~\ref{prop:quadratic-summability} shows
is actually needed.  It does not rule out a construction that avoids
activation altogether and instead reads $d_q(\sigma)$ off a quantity that
already exists in the cascade as an error term, rather than as an amplified
response.
\end{remark}

\begin{question}[Natural-residual injection]
\label{q:natural-residual}
The cascade of \cite{AlpogeBuckmasterCoiculescu2026} already produces, at
each level, force increments and structured defects $G_{q+1}$, $E_j$, $Q_j$
(its eqs.~(134)--(135), (154)) that are not designed but forced by the
mismatch between the truncated parametrix and the exact equation.  Is there
a choice of the fixed profiles $f,\chi,W$, together with a rule identifying
the activation times $s_{q+1}$ with the squarefree times $s_q$ of
Section~\ref{sec:squarefree-pairs}, for which one of these already-occurring
residual quantities, read at $(0,1)$, equals $\lambda_\sigma\,c_qs_q^{-\sigma}+e_q$
for a fixed nonzero constant $\lambda_\sigma$ independent of $q$ and
$\sigma$, with $\sum_q|e_q|<\infty$, without choosing any function of the
construction using the value $c_q$ itself?
\end{question}

Question~\ref{q:natural-residual} is \textup{[OPEN, speculative]} and is, on
present evidence, the only candidate not excluded by
Propositions~\ref{prop:sign-obstruction} and \ref{prop:size-obstruction}:
residual terms are not subject to the activation-gain lower bound of
Lemma~4.9, since they are the leftover error of an approximation rather
than the leading amplified response.  Whether any such residual is large
enough to be read out at all, and small enough to match
Proposition~\ref{prop:quadratic-summability}'s size, is not addressed by any
source consulted here.  Answering Question~\ref{q:natural-residual} in the
negative --- showing every such residual is either too small to read
(unconditionally smaller than the estimates already used to prove force
smoothness) or itself tied to an amplifying step, reintroducing
Proposition~\ref{prop:size-obstruction} --- would be as informative as
answering it in the affirmative.

\begin{warning}[What a completed insertion would and would not show about RH]
\label{warn:no-progress-without-independent-bound}
Suppose Question~\ref{q:parent-insertion} were answered completely: a
genuine, forced, jointly smooth family of Euler or IPM solutions were built
whose real deformation ratio equals $\exp(\delta_Q(\sigma))$ exactly, for
every $\sigma>3/2$.  By Theorem~\ref{thm:conditional-PDE-transfer}, this
would upgrade statement~(iv) of Theorem~\ref{thm:CMZ-RH-equivalence} from a
toy-ODE statement to a genuine PDE statement, closing
Question~\ref{q:paired-reunion}.  It would \emph{not} decide the Riemann
hypothesis, and would not bring a decision closer, for a specific reason:
the equivalence of RH with convergence of $\delta_Q(\sigma)$ is already
proved unconditionally in Theorem~\ref{thm:CMZ-RH-equivalence} by
Littlewood's criterion alone, with no reference to any PDE.  A PDE
realisation re-encodes $\delta_Q(\sigma)$; it supplies no independent
information about $M(x)$.  Whether $\delta_Q(\sigma)$ converges would remain
exactly as open after the realisation as before it.  Propositions
\ref{prop:sign-obstruction} and \ref{prop:size-obstruction} illustrate the
same point concretely: the two most natural native candidates fail not from
absence of a proof but from outright conflict with already-proved lemmas of
the cited sources, which is itself evidence that the activated-packet
mechanism is the wrong tool for producing a small, freely signed quantity ---
Tao's own description of the method (blog post of 7 September 2026) is that
it is designed to locate an ``exploitable instability'' in a background
flow, which is definitionally the opposite goal from bounding how much
instability a construction can produce.

The only way a construction of this kind could contribute to deciding RH is
if the completed PDE side supplied an inequality that is not derived by
first estimating $c_q$ or $U(q)$ arithmetically and then importing the
estimate into the packets, but is instead a bound holding for every
admissible family of auxiliary packets of the type in
\eqref{eq:parent-insertion-target}, proved by genuinely fluid-dynamical
means (an energy identity, a monotonicity formula, a maximum principle for
the transported angle or amplitude), and that, when specialised to the
ordinal-prime packets, gives a bound on $M(x)$ at least as strong as
$O_\eta(x^{1/2+\eta})$.  No such universal bound exists in the literature
consulted here: the multiscale method of
\cite{CordobaMartinezZoroa2023,CordobaMartinezZoroaZheng2023,AlpogeBuckmasterCoiculescu2026,AlpogeBuckmaster2026}
proves \emph{existence} of a blow-up for a chosen forcing, which is the
opposite kind of statement from a \emph{universal} bound holding for every
admissible forcing.  \textup{[OPEN, speculative]}: whether any
fluid-mechanical monotonicity of this kind exists for the CMZ mechanism is
not addressed by the present manuscript.  Absent such a bound, every result
in this section, however completed, is a physical model of the Mertens
problem and not a route to its resolution.
\end{warning}

\begin{question}[Parent-insertion problem]
\label{q:parent-insertion}
Can auxiliary packets, attached to the cascade of
\cite{AlpogeBuckmasterCoiculescu2026} at the activation times of the ordinal
pairs, be constructed so that \eqref{eq:parent-insertion-target} holds with
$\sum_q|E_q|\,(s_q^{-\sigma}-s_{q+1}^{-\sigma})<\infty$, while the resulting
force increments remain summable in the sense used in
Proposition~\ref{prop:exact-CMZ-readout}?
\end{question}

Question~\ref{q:parent-insertion} is \textup{[OPEN]} and replaces
Question~\ref{q:paired-reunion} as the concrete remaining problem; per
Warning~\ref{warn:no-progress-without-independent-bound}, closing it would
settle Question~\ref{q:paired-reunion} without settling RH.
Propositions~\ref{prop:sign-obstruction} and~\ref{prop:size-obstruction}
show that the natural ways of specialising
Question~\ref{q:parent-insertion} to the activated-packet mechanism fail
outright; Question~\ref{q:natural-residual} is the specific remaining
instance not excluded by either.

\section{Boundary, convergence, and interchange audit}
\label{sec:audit}

\begin{proposition}[Complete audit of the logical dependencies]
\label{prop:dependency-audit}
The proof has the following dependency chain:
\begin{equation}
(S,SS)\longrightarrow U(q)
\longrightarrow\delta_Q(\sigma)
\longrightarrow\mathscr C_Q(\sigma)
\longrightarrow\text{tail stability}
\longrightarrow\RH.
\label{eq:dependency-chain}
\end{equation}
The first three arrows are finite identities.  The fourth is the exact scalar
ODE solution.  The final arrow uses the original-order convergence theorem
and the Littlewood criterion.  No step before tail stability assumes RH.
\end{proposition}

\begin{proof}
The first arrow is Proposition~\ref{prop:U-open-block}; the second is
Lemma~\ref{lem:finite-Abel}; the third is Proposition
\ref{prop:exact-CMZ-readout}; and the last two arrows are Theorem
\ref{thm:CMZ-RH-equivalence}.
\end{proof}

The analytic and finite operations used above have the following status.
\begin{enumerate}[label=\textup{A\arabic*.},leftmargin=*]
\item All rearrangements defining $A+B+C+\delta$ take place first in
$\Re s>2$, where coefficient bounds give absolute convergence.
\item Formula \eqref{eq:finite-Abel-identity} is finite and valid for every
complex $s$; it requires no convergence theorem.
\item The physical CMZ clock uses real $\sigma>0$, ensuring positivity of all
interval lengths.
\item The unfinished-pair boundary is retained as $U(Q)s_Q^{-\sigma}$.
Under RH it vanishes by \eqref{eq:RH-W-bound}; in the converse direction it is
part of $\delta_Q$ and is not separately discarded.
\item The integral defining $C$ is holomorphic in $\Re s>3/2$ by the explicit
majorant \eqref{eq:C-majorant}.  No continuation of $\delta$ is used in that
argument.
\item Singularities of $B$ are not promoted to singularities of $\alpha$.
The analytic criterion concerns the proved continuation of
$A+B+C=\alpha-\delta$.
\item The CMZ literature estimates concern prescribed PDE layers.  They do not
establish Definition~\ref{def:admissible-readout}; that definition is an
additional hypothesis in Theorem~\ref{thm:conditional-PDE-transfer}.
\end{enumerate}

\begin{remark}[A failed shortcut]
\label{rem:failed-shortcut}
It is insufficient to prove merely that each finite arithmetic profile can be
inserted into a CMZ layer after choosing an arbitrarily large carrier scale.
Such a choice may make the PDE residual small while the normalised logarithmic
deformation still grows with $Q$.  The uniform, unattenuated identity
\eqref{eq:PDE-readout-error} and the tail estimate
\eqref{eq:readout-error-Cauchy} are the properties that discriminate the RH
problem from finite gluing flexibility.
\end{remark}

\section{Conclusion and research route}
\label{sec:conclusion}

The ordinal-prime transform produces an exact open-pair strain $U(q)$ that is
readable from $SS$.  Finite Abel summation turns this strain into the residual
Dirichlet polynomial $\delta_Q$.  Splitting $U$ into positive and negative
parts permits two positive logarithmic-stretching channels, and their relative
deformation equals $\exp(\delta_Q)$.  Theorem
\ref{thm:CMZ-RH-equivalence} is therefore an exact dynamical reformulation of
RH, with the physical clock confined to the fixed interval $[0,1]$.

The unresolved PDE route is specific: realise the paired arithmetic strains
by genuine transported CMZ packets and prove a Cauchy tail estimate for their
relative deformation without allowing carrier frequencies to attenuate the
arithmetic coefficient, now stated as the parent-insertion problem,
Question~\ref{q:parent-insertion}.  This is an \textup{[OPEN]} research
proposal.  It is not supplied by the four fluid papers cited here, and no
claim of novelty relative to all literature is made.  Three more direct
attempts at an insertion (Section~\ref{subsec:three-attempts}) were checked
against the cited theorems and each falls outside what those theorems prove.
Whatever the fate of Question~\ref{q:parent-insertion}, it does not bear on
RH itself: this is stated precisely in
Warning~\ref{warn:no-progress-without-independent-bound}, and is the single
most important qualification in this manuscript.

The route count is as follows: four \textup{[LIT]} multiscale inputs, one
\textup{[ADAPT]} two-channel stretching model, two principal \textup{[DED]}
bridges (finite readout and RH equivalence), one \textup{[DED]}
generalisation of an existing lemma (Lemma~\ref{lem:generalised-layer-bound}),
one \textup{[DED]} unconditional summability bound
(Proposition~\ref{prop:quadratic-summability}), four examined insertion
attempts of which two are excluded outright by direct conflict with
already-proved lemmas
(Propositions~\ref{prop:sign-obstruction},~\ref{prop:size-obstruction}) and
two remain unproved extensions, one \textup{[OPEN, speculative]} candidate
not excluded by either obstruction
(Question~\ref{q:natural-residual}), and one explicitly \textup{[OPEN]} PDE
realisation problem whose resolution would not, by itself, decide RH.

\section*{Notebook History}
\label{sec:notebook-history}

Version 1.0: definitions of $\zeta$, $\mu$, $\lambda$, and $\Lambda$;
fundamental properties of $S$, $SS$, and $\alpha$; audited CMZ proof summary;
finite-boundary readout; full CMZ deformation equivalence; analytic
singularity audit; and conditional genuine-PDE transfer.

Version 2.0: examined three candidate insertions of the ordinal-prime data
into a genuine forced Euler/IPM tower (amplitude rescaling, continuous
exponent encoding, parent-gradient insertion), and showed each is either an
unproved extension of the cited theorems or an encoding that does not
constrain the arithmetic input; proved a bounded-ratio generalisation of the
layered self-interaction bound of C\'ordoba, Mart\'inez-Zoroa, and Zheng;
replaced the paired-reunion question by the parent-insertion problem; and
recorded, as a standalone warning, the precise sense in which completing any
of these routes would not by itself decide the Riemann hypothesis.

Version 3.0: proved that reading the ordinal-prime data off a single
tower's amplitude-gain ratio is impossible, not merely unproved, by direct
conflict with the amplitude-gain lemma of Alp\"oge, Buckmaster, and
Coiculescu; proved an unconditional quadratic-summability bound for the
ordinal-prime target sequence; used it to show a two-tower comparison of
activated packets is excluded for the same structural reason, restated
precisely; and isolated residual, non-amplified terms of the construction
as the one candidate excluded by neither obstruction, posed as the
natural-residual injection question.

\begin{thebibliography}{99}

\bibitem{Apostol1976}
T. M. Apostol,
\emph{Introduction to Analytic Number Theory},
Undergraduate Texts in Mathematics, Springer, New York, 1976.
DOI: 10.1007/978-1-4757-5579-4.

\bibitem{Littlewood1912}
J. E. Littlewood,
Quelques cons\'equences de l'hypoth\`ese que la fonction $\zeta(s)$ de
Riemann n'a pas de z\'eros dans le demi-plan $\Re s>1/2$,
\emph{C. R. Acad. Sci. Paris} \textbf{154} (1912), 263--266.

\bibitem{Titchmarsh1986}
E. C. Titchmarsh,
\emph{The Theory of the Riemann Zeta-Function},
2nd ed., revised by D. R. Heath-Brown, Clarendon Press, Oxford, 1986.

\bibitem{Ng2004}
N. Ng,
The distribution of the summatory function of the M\"obius function,
\emph{Proc. London Math. Soc.} (3) \textbf{89} (2004), 361--389.
DOI: 10.1112/S0024611504014741; arXiv:math/0310381.

\bibitem{CordobaMartinezZoroa2023}
D. C\'ordoba and L. Mart\'inez-Zoroa,
Blow-up for the incompressible 3D-Euler equations with uniform
$C^{1,1/2-\varepsilon}\cap L^2$ force,
preprint, 2023.
arXiv:2309.08495.

\bibitem{CordobaMartinezZoroaZheng2023}
D. C\'ordoba, L. Mart\'inez-Zoroa, and F. Zheng,
Finite time singularities to the 3D incompressible Euler equations for
solutions in $C^\infty(\mathbb R^3\setminus\{0\})\cap C^{1,\alpha}\cap L^2$,
preprint, 2023.
arXiv:2308.12197.

\bibitem{AlpogeBuckmasterCoiculescu2026}
L. Alp\"oge, T. Buckmaster, and M. P. Coiculescu,
Extending the C\'ordoba--Mart\'inez-Zoroa IPM blow-up to uniformly space-time
smooth forcing,
manuscript, 2026.

\bibitem{AlpogeBuckmaster2026}
L. Alp\"oge and T. Buckmaster,
Blowup for the Boussinesq equations with smooth forcing,
manuscript, 2026.

\end{thebibliography}

\end{document}
